\chapter{Differential Geometry of Surfaces}\label{ch:diffgeo}

Normals, curvatures and geodesic distances are the vocabulary of geometry processing: smoothing
reduces curvature, simplification preserves it, remeshing aligns elements with principal
directions, and shape descriptors summarise it. This chapter reviews the differential geometry of
curves and surfaces and shows how each quantity is computed on a triangle mesh. Classical
references are \citet{doCarmo1976} for the smooth theory and \citet{Meyer2003,Crane2018} for the
discrete operators. We use the analyst's Laplacian $\Lap=\Div\grad$ in this chapter.

\section{Curves}

Let $\gamma:[a,b]\to\R^3$ be a regular curve, $\gamma'(t)\neq\bzero$. Its \emph{arc length}
$s(t)=\int_a^t\norm{\gamma'(\tau)}\dd\tau$ allows a reparameterization with unit speed, in which
the unit tangent is $\bt(s)=\gamma'(s)$. The \emph{curvature} measures how fast the tangent turns,
\begin{equation}
\kappa(s)=\norm{\bt'(s)},\qquad \bt'(s)=\kappa(s)\,\bn(s),
\end{equation}
where $\bn$ is the principal normal. The \emph{osculating circle} at $\gamma(s)$ is the circle
tangent to the curve with radius $1/\kappa$ and centre $\gamma(s)+\bn(s)/\kappa(s)$: it is the
circle that best approximates the curve to second order. For a planar curve the curvature can be
given a sign, positive when the curve turns counter-clockwise: if $\theta(s)$ is the angle of the
tangent, then $\kappa=\dd\theta/\dd s$.

\paragraph{Discrete curvature of a polyline.} For a polygon with vertices $\bp_i$, the tangent is
constant on each edge and turns by the exterior angle $\theta_i$ at each vertex. Curvature is
therefore concentrated at the vertices as a Dirac mass, and the natural discrete quantity is
the \emph{integrated} curvature $\theta_i=\int\kappa\dd s$ over the dual cell of the vertex. Dividing
by the length of that cell, $\ell_i=\tfrac12(\norm{\bp_{i+1}-\bp_i}+\norm{\bp_i-\bp_{i-1}})$, gives a
pointwise estimate $\kappa_i\approx\theta_i/\ell_i$ that converges as the polygon is refined
(\cref{fig:curve}). A useful exact identity follows: for a closed planar polygon the turning angles
add up to $2\pi$ times the turning number, the discrete version of $\oint\kappa\dd s=2\pi k$.

\fig[0.85\textwidth]{ch04_curve.pdf}{Left: a closed planar curve with osculating circles of
radius $1/\abs{\kappa}$ at four points; the arrows are the tangent and the normal. Right: exact
signed curvature against arc length, and the discrete estimate (turning angle divided by the dual
length) on polygons with $n=24$ and $n=60$ vertices.}{fig:curve}

\section{Surfaces}

\subsection{Tangent plane and normal}

Let $\bx(u,v)$ be a regular parameterization of a surface patch. The partial derivatives $\bx_u$
and $\bx_v$ span the tangent plane, and the unit normal is
\begin{equation}
\bn=\frac{\bx_u\times\bx_v}{\norm{\bx_u\times\bx_v}} .
\end{equation}

\subsection{First fundamental form}

A tangent vector $\bw=a\bx_u+b\bx_v$ has squared length
\begin{equation}
\norm{\bw}^2=\begin{pmatrix}a&b\end{pmatrix}\mathbf{I}\begin{pmatrix}a\\b\end{pmatrix},\qquad
\mathbf{I}=\begin{pmatrix}E&F\\F&G\end{pmatrix}=\begin{pmatrix}\bx_u\cdot\bx_u&\bx_u\cdot\bx_v\\
\bx_u\cdot\bx_v&\bx_v\cdot\bx_v\end{pmatrix}.
\end{equation}
The first fundamental form $\mathbf I$ (the \emph{metric}) determines lengths of curves, angles
between tangent vectors and areas, with area element $\dd A=\sqrt{EG-F^2}\dd u\dd v$. Quantities
that depend only on $\mathbf I$ are called \emph{intrinsic}: they can be measured by an inhabitant
of the surface without looking at the ambient space, and they do not change when the surface is
bent without stretching (an \emph{isometric} deformation). Geodesic distances, the Gaussian
curvature and the Laplace--Beltrami operator are intrinsic; this is the basis of the
pose-invariant descriptors of \cref{ch:spectral}.

\subsection{Second fundamental form and the shape operator}

The second fundamental form measures how the surface bends away from its tangent plane:
\begin{equation}
\mathbf{II}=\begin{pmatrix}e&f\\f&g\end{pmatrix}=\begin{pmatrix}\bx_{uu}\cdot\bn&\bx_{uv}\cdot\bn\\
\bx_{uv}\cdot\bn&\bx_{vv}\cdot\bn\end{pmatrix}.
\end{equation}
Intersecting the surface with the plane spanned by $\bn$ and a unit tangent direction $\bw$ gives a
curve whose curvature is the \emph{normal curvature}
\begin{equation}
\kappa_n(\bw)=\frac{\mathbf{II}(\bw,\bw)}{\mathbf{I}(\bw,\bw)} .
\end{equation}
The \emph{shape operator} (Weingarten map) $\mS=\mathbf{I}^{-1}\mathbf{II}$ is self-adjoint with
respect to the metric. Its eigenvalues $\kappa_1\ge\kappa_2$ are the \emph{principal curvatures},
the maximum and minimum of $\kappa_n$, and its eigenvectors are the orthogonal \emph{principal
directions} $\bt_1,\bt_2$. Writing $\bw=\cos\varphi\,\bt_1+\sin\varphi\,\bt_2$ gives Euler's
theorem
\begin{equation}
\kappa_n(\varphi)=\kappa_1\cos^2\varphi+\kappa_2\sin^2\varphi .
\end{equation}

\begin{definition}[Mean and Gaussian curvature]
$H=\tfrac12(\kappa_1+\kappa_2)=\tfrac12\tr\mS$ and
$K=\kappa_1\kappa_2=\det\mS=\dfrac{eg-f^2}{EG-F^2}$.
\end{definition}

Conversely $\kappa_{1,2}=H\pm\sqrt{H^2-K}$. The sign of $K$ classifies points: \emph{elliptic}
($K>0$, the surface lies locally on one side of the tangent plane, like a sphere or the top of a
hill), \emph{hyperbolic} ($K<0$, a saddle), \emph{parabolic} ($K=0$ with one non-zero principal
curvature, like a cylinder) and \emph{planar} ($\kappa_1=\kappa_2=0$). Points with
$\kappa_1=\kappa_2$ are \emph{umbilics}; there the principal directions are undefined, which
creates singularities in direction fields (\cref{ch:remesh}).

\begin{example}{Sphere, cylinder and torus}{curvatures}
A sphere of radius $r$ with outward normal has $\kappa_1=\kappa_2=1/r$ (with the sign convention
that convex surfaces have positive curvature), so $H=1/r$ and $K=1/r^2$; every point is umbilic. A
cylinder of radius $r$ has $\kappa_1=1/r$ around it and $\kappa_2=0$ along it, so $H=1/(2r)$ and
$K=0$. For the torus $\bx(u,v)=((R+r\cos v)\cos u,(R+r\cos v)\sin u,r\sin v)$ the principal
curvatures are $1/r$ (around the tube) and $\cos v/(R+r\cos v)$ (around the axis), hence
\begin{equation}
K=\frac{\cos v}{r(R+r\cos v)},\qquad H=\frac{R+2r\cos v}{2r(R+r\cos v)} .
\end{equation}
The outer half ($\cos v>0$) is elliptic, the inner half hyperbolic, and the top and bottom circles
parabolic. These exact values are the ground truth of \cref{fig:torus-curv}.
\end{example}

\subsection{Two global theorems}

\begin{theorem}[Theorema Egregium, Gauss]
The Gaussian curvature depends only on the first fundamental form, so it is invariant under
isometric deformations.
\end{theorem}
A sheet of paper ($K=0$) can be rolled into a cylinder ($K=0$) but never wrapped around a sphere
($K>0$) without stretching. The mean curvature, in contrast, is extrinsic: rolling the sheet
changes $H$ from $0$ to $1/(2r)$.

\begin{theorem}[Gauss--Bonnet]
For a closed orientable surface $S$, $\displaystyle\int_SK\dd A=2\pi\chi(S)=2\pi(2-2g)$. For a
surface with boundary, $\displaystyle\int_SK\dd A+\oint_{\partial S}k_g\dd s=2\pi\chi(S)$, where
$k_g$ is the geodesic curvature of the boundary.
\end{theorem}
The total curvature is a topological invariant: denting a sphere creates negative curvature that is
exactly compensated by positive curvature elsewhere.

\subsection{Mean curvature and the Laplace--Beltrami operator}

On a surface, the Laplace--Beltrami operator $\Lap_S=\Div_S\grad_S$ generalises the Laplacian. Its
action on the coordinate functions has a remarkable geometric meaning:
\begin{equation}
\Lap_S\bx=-2H\bn ,
\label{eq:mean-curvature-normal}
\end{equation}
with $\bn$ the outward normal and $H>0$ on convex regions. On the unit sphere the coordinates are
spherical harmonics of degree one, $\Lap_S\bx=-2\bx$, in agreement with $H=1$. Identity
\eqref{eq:mean-curvature-normal} explains why diffusing the coordinates smooths a surface: the
flow $\partial_t\bx=\Lap_S\bx$ is the \emph{mean curvature flow}, which moves every point against
the normal with speed $2H$ and flattens bumps (\cref{ch:smooth}). It also identifies minimal
surfaces ($H=0$) as surfaces whose coordinates are harmonic functions.

\section{Discrete normals}

The normal of a triangle $(\bx_a,\bx_b,\bx_c)$ is
$\bn_f=(\bx_b-\bx_a)\times(\bx_c-\bx_a)/\norm{\cdot}$, and the norm of the unnormalised cross
product is twice the area. A mesh has no well-defined normal at the vertices, so vertex normals
are weighted averages of the normals of the incident faces,
\begin{equation}
\bn_i=\frac{\sum_{f\ni i}w_f\bn_f}{\norm{\sum_{f\ni i}w_f\bn_f}},
\end{equation}
with common choices \citep{Jin2005}:
\begin{itemize}
  \item \textbf{uniform} weights $w_f=1$ (cheap, but depends on how the one-ring is triangulated);
  \item \textbf{area} weights $w_f=A_f$, obtained for free by summing unnormalised cross products;
  \item \textbf{angle} weights $w_f=\theta_f$, the interior angle of $f$ at vertex $i$
  \citep{Thurmer1998}, which do not change when a face is subdivided;
  \item the weights of \citet{Max1999}, $w_f=\sin\theta_f/(\norm{\be_1}\norm{\be_2})$, exact for
  vertices on a sphere.
\end{itemize}
On a smooth surface sampled by a reasonable triangulation all of them converge to the true normal
as the mesh is refined (\cref{fig:normals}). The differences matter on coarse, irregular or
anisotropic meshes and near sharp features, where none of the averages is meaningful and the
normal should be duplicated per face group.

\fign{ch04_normals.pdf}{Convergence of vertex normals on irregularly triangulated ellipsoids: the
95th percentile of the angle to the exact normal against the mean edge length $h$. On this smooth
surface all three weightings converge at about second order, with angle weights slightly more
accurate on the finest meshes.}{fig:normals}

\section{Discrete curvature}

\subsection{Gaussian curvature: the angle defect}

Flatten the triangles around an interior vertex $i$ into the plane: their angles $\theta_j$ add
up to $2\pi$ only if the one-ring is flat. The \emph{angle defect}
\begin{equation}
\Omega_i=2\pi-\sum_{j}\theta_j \qquad(\pi-\sum_j\theta_j\ \text{at boundary vertices})
\end{equation}
is the integrated Gaussian curvature of the dual cell, and $K_i=\Omega_i/A_i$, with $A_i$ the area
associated to the vertex, is the pointwise estimate. A cone vertex such as the apex of a pyramid
has positive defect, a saddle vertex has more than $2\pi$ of angle around it and negative defect.

\begin{proposition}[Discrete Gauss--Bonnet]
For any closed triangle mesh, $\sum_i\Omega_i=2\pi\chi$.
\end{proposition}
\begin{proof}
The angles of all triangles add up to $\pi F$, so $\sum_i\Omega_i=2\pi V-\pi F$. Since $2E=3F$,
$V-E+F=V-\tfrac12F$, hence $\sum_i\Omega_i=2\pi(V-E+F)=2\pi\chi$.
\end{proof}
The statement is exact, not an approximation: on the torus of \cref{fig:torus-curv} the sum of
angle defects is $-6.7\times10^{-13}$, zero up to floating-point error.

\subsection{Vertex areas}

The area $A_i$ assigned to a vertex should partition the surface. The \emph{barycentric} area is
one third of the area of the incident triangles. The \emph{Voronoi} area, the part of each triangle
closer to $i$ than to the other corners, is
\begin{equation}
A_i^{\mathrm{Vor}}=\frac18\sum_{j\in\mathcal N(i)}(\cot\alpha_{ij}+\cot\beta_{ij})\norm{\bx_i-\bx_j}^2,
\end{equation}
where $\alpha_{ij},\beta_{ij}$ are the angles opposite the edge $ij$. It fails on obtuse
triangles, where the circumcentre falls outside the triangle, so \citet{Meyer2003} define the
\emph{mixed} area: the Voronoi area for non-obtuse triangles and $A_f/2$ or $A_f/4$ for obtuse
triangles, depending on whether the obtuse angle is at $i$.

\subsection{Mean curvature: the cotangent formula}

Discretising \eqref{eq:mean-curvature-normal} with linear finite elements (\cref{ch:smooth}) gives
the \emph{cotangent formula} \citep{Pinkall1993,Meyer2003}
\begin{equation}
(\Lap\bx)_i=\frac{1}{2A_i}\sum_{j\in\mathcal N(i)}(\cot\alpha_{ij}+\cot\beta_{ij})(\bx_j-\bx_i)
\approx-2H_i\bn_i ,
\label{eq:cotan}
\end{equation}
so $\abs{H_i}=\tfrac12\norm{(\Lap\bx)_i}$ and its sign is the sign of $-(\Lap\bx)_i\cdot\bn_i$. The
vector $(\Lap\bx)_i$ is also a normal estimate, the \emph{mean curvature normal}, but it vanishes
on flat regions. With $H$ and $K$ the principal curvatures follow from $\kappa_{1,2}=H\pm\sqrt{H^2-K}$
(clamping $H^2-K$ at zero, since the two estimates are computed independently).

\fig{ch04_torus_curvature.png}{Discrete curvatures on a torus with $R=1$, $r=0.35$. The angle
defect reproduces the exact Gaussian curvature (positive outside, negative inside) and the
cotangent formula the exact mean curvature; the scatter plots against the closed-form values of
\cref{ex:curvatures} fall on the diagonal.}{fig:torus-curv}

\fig[0.75\textwidth]{ch04_model_curvature.png}{Mean and Gaussian curvature of the bunny. Red is
positive and blue negative. Ridges such as the rim of the ears have large $H$; the neck and the
junctions between parts are hyperbolic ($K<0$). The colour ranges are clipped at the 1st and 99th
percentiles because curvature is very noisy on scanned meshes.}{fig:model-curv}

\subsection{The curvature tensor and principal directions}

Principal directions require the full shape operator, a $2\times2$ symmetric matrix in a local
tangent frame. Three families of estimators are common:
\begin{itemize}
  \item \textbf{Fitting}: express the neighbours of a vertex in a frame $(\be_1,\be_2,\bn)$ and fit
  a height function $z=\tfrac12(ax^2+2bxy+cy^2)$ by least squares; then
  $\mS\approx\begin{psmallmatrix}a&b\\b&c\end{psmallmatrix}$. Higher-order polynomial fits
  (osculating jets) are more accurate on noisy data \citep{Cazals2005}.
  \item \textbf{Normal variation}: estimate $\mS$ per face from the change of the vertex normals
  along the edges and average per vertex \citep{Rusinkiewicz2004}.
  \item \textbf{Edge-based tensors}: sum, over the edges in a neighbourhood, the dihedral angle
  times the outer product of the edge direction \citep{Taubin1995b,Cohen2003}; the result is a
  $3\times3$ tensor whose two largest eigenvectors are the principal directions (swapped).
\end{itemize}
\Cref{fig:principal} shows principal directions computed by quadric fitting. On the ellipsoid
they form the familiar lines-of-curvature pattern with four umbilics; on the torus $\bt_1$ is
everywhere tangent to the meridians.

\fig{ch04_principal.png}{Principal directions $\bt_1$ estimated by local quadric fitting. Left:
an ellipsoid coloured by $K$. Centre: a torus coloured by $\kappa_2$. Right: the cow Spot coloured
by the anisotropy $\abs{\kappa_1-\kappa_2}$, which is large on the legs and the horns, where the
principal directions are well defined.}{fig:principal}

\section{Geodesic distance}

The geodesic distance between two points is the length of the shortest curve on the surface that
joins them. It is intrinsic and it is the ground truth of many experiments in
\cref{ch:deform,ch:spectral,ch:fmaps}. Three ways of computing it on a mesh are:
\begin{itemize}
  \item \textbf{Dijkstra} on the edge graph: simple and fast, but paths are restricted to edges, so
  distances are overestimated by a few percent (metrication error) that does not vanish under
  refinement.
  \item \textbf{Fast marching} \citep{Kimmel1998}: solves the eikonal equation
  $\norm{\grad u}=1$ by propagating a front through the triangles, with first-order accuracy.
  \item \textbf{The heat method} \citep{Crane2013}: (i) diffuse heat from the source for a short
  time $t\approx h^2$ by solving $(\mM-t\mC)\bu=\bdelta$; (ii) normalise the gradient,
  $\bm X=-\grad u/\norm{\grad u}$; (iii) solve the Poisson equation $\Lap\phi=\Div\bm X$. Only
  sparse linear solves with the cotangent matrix are needed, and the factorisations can be reused
  for every source. The figures of these notes use the heat method with $t=4h^2$.
\end{itemize}
The \emph{vector heat method} \citep{Sharp2019} extends the same idea to the transport of tangent
vectors.

\begin{keypoints}
\begin{itemize}[leftmargin=*]
  \item The first fundamental form gives lengths, angles and areas; the second gives normal
  curvatures. The shape operator $\mI^{-1}\mathbf{II}$ has eigenvalues $\kappa_1,\kappa_2$.
  \item $K$ is intrinsic (Theorema Egregium) and integrates to $2\pi\chi$ (Gauss--Bonnet); $H$ is
  extrinsic and appears in $\Lap_S\bx=-2H\bn$.
  \item On meshes: angle defect for $K$ (with an exact discrete Gauss--Bonnet theorem), cotangent
  formula for $H\bn$, mixed Voronoi areas for normalisation, fitting for principal directions.
  \item Geodesic distances can be computed by Dijkstra (biased), fast marching or the heat method.
\end{itemize}
\end{keypoints}

\section*{Exercises}
\begin{exercise}
Compute $\mathbf I$, $\mathbf{II}$, $H$ and $K$ for the paraboloid $z=\tfrac12(ax^2+by^2)$ at the
origin and interpret the result.
\end{exercise}
\begin{exercise}
Integrate the Gaussian curvature of the torus over $u,v\in[0,2\pi)$ using $\dd A=r(R+r\cos v)\dd u\dd v$
and verify Gauss--Bonnet.
\end{exercise}
\begin{exercise}
Show that the angle defect of the apex of a regular pyramid with $k$ lateral faces tends to zero as
the pyramid is flattened and to $2\pi$ as it becomes a needle.
\end{exercise}
\begin{exercise}
Prove that for an equilateral triangle mesh of the plane the cotangent weights are all
$2\cot(\pi/3)=2/\sqrt3$, so the cotangent and uniform Laplacians coincide up to a constant.
\end{exercise}
